%% ════════════════════════════════════════════════════════════════════════════
\section{Introduction}
\label{sec:intro}
%% ════════════════════════════════════════════════════════════════════════════

\subsection{Context and Motivation}
\label{subsec:context}

The growing adoption of Internet of Things (IoT) technology in education
has opened new possibilities for automating routine classroom tasks and
improving the quality of the learning environment.  Despite this trend,
many universities---including SMU\,--\,Mediterranean Institute of
Technology---continue to rely on manual processes for attendance
management and environmental control.  Professors spend up to 15~minutes
per lecture calling roll, consuming valuable instructional time.
Meanwhile, classroom HVAC systems and lighting are operated without any
sensor feedback, leaving occupant comfort dependent on manual judgment.
Early identification of students at risk of academic failure due to poor
attendance remains largely reactive, with intervention arriving after the
damage is already done.

The convergence of affordable edge hardware (ESP32, Raspberry Pi~4B),
open-source ML frameworks (DeepFace, Ollama), and lightweight message
protocols (MQTT) makes it feasible to deploy a fully autonomous
classroom management platform at minimal cost, running entirely on the
university's local area network with no cloud dependency.

\subsection{Problem Statement}
\label{subsec:problem}

Three concrete pain points motivate this work:

\begin{enumerate}
  \item \textbf{Attendance inefficiency.}  Manual roll-call consumes up
    to 15~minutes per session---time that reduces effective teaching time
    by 15--20\% for a 50-minute lecture.  Records are entered manually
    into the Moodle LMS, creating duplication of effort and data-entry
    errors.

  \item \textbf{Unmonitored environment.}  No real-time sensor data is
    available to inform HVAC or lighting decisions.  Temperature, humidity,
    CO$_2$/VOC levels, and occupancy are unknown to both professors and
    facility managers.

  \item \textbf{No early warning for at-risk students.}  Advising staff
    receive no automated signal when a student's attendance rate begins to
    decline.  Intervention is reactive rather than preventive, reducing
    its effectiveness.
\end{enumerate}

\subsection{Project Objectives}
\label{subsec:objectives}

The general objective of this project is to design and implement an
IoT-based smart classroom management system that automates student
attendance via camera-based face recognition, monitors classroom
environmental conditions in real time, enables automated and manual
control of AC and lighting, provides professors with a live web
dashboard, synchronises attendance to the Moodle LMS, and identifies
at-risk students and attendance trends using a locally hosted LLM---all
running on a Raspberry Pi~4B without cloud dependency.

The six specific objectives are:
\begin{enumerate}
  \item Automate student attendance recording using face recognition.
  \item Monitor classroom environmental conditions (temperature, humidity,
    air quality, sound) in real time.
  \item Enable automated and manual control of AC and lighting via relay
    actuation.
  \item Provide professors with a live web dashboard showing attendance,
    sensor data, and alerts.
  \item Synchronise attendance records to Moodle automatically on session
    end.
  \item Identify at-risk students and forecast attendance trends using AI
    analytics.
\end{enumerate}

\subsection{Scope and Limitations}
\label{subsec:scope}

The system is scoped to a single classroom room (\texttt{room1}) on the
university intranet.  All services run locally on the Raspberry Pi~4B;
there is no internet dependency during normal operation.  The academic
semester timeline constrained the project to a proof-of-concept
deployment with physical hardware; production hardening (HTTPS/WSS,
multi-room support) is identified as future work.

Hardware budget constraints led to the selection of the MQ-135 gas
sensor, whose output is uncalibrated and used comparatively rather than
as an absolute CO$_2$ reference.  The Raspberry Pi~4B's CPU limits
face-recognition throughput to 2~fps and phi3-mini inference to
10--30~seconds per student.

\subsection{Report Structure}
\label{subsec:structure}

Section~\ref{sec:literature} reviews related work on IoT-based attendance
systems, face recognition, and edge AI.
Section~\ref{sec:requirements} presents stakeholder analysis and
functional and non-functional requirements.
Section~\ref{sec:architecture} describes the full system architecture
across all subsystem layers.
Section~\ref{sec:implementation} details hardware assembly, firmware, and
software implementation.
Section~\ref{sec:ai} covers the AI analytics pipelines.
Section~\ref{sec:security} analyses security and privacy considerations.
Section~\ref{sec:testing} documents the testing and validation approach.
Section~\ref{sec:results} presents results and key lessons learned.
Section~\ref{sec:conclusion} concludes with a summary of contributions
and directions for future work.
